\begin{textsample}{NEG Dim 8 – global\_north\_2024\_04 – Score -2.00 – global\_north\_2024\_04/107573477.txt}  \label{ex:f8_neg_130}
Qatar ' s warning that it was reassessing its role as a mediator between Israel and Hamas has raised concerns about the prospects for a ceasefire and the return of hostages.
Since Hamas ' s deadly October 7 attack on Israel sparked devastating retaliation against Gaza, Qatar has been an important \textbf{conduit} to the Palestinian militant group, whose political office is in Doha.
The gas-rich emirate, which also hosts the region ' s largest US military base, successfully brokered a \textbf{week-long} break in fighting in late November when scores of Israeli and foreign hostages were released.
But with months of further negotiations failing to win a truce, and with Qatar facing criticism notably from Israel, Prime Minister Mohammed bin Abdulrahman Al Thani said on Wednesday it was time for a"complete re-evaluation"of its role. - What prompted Qatar ' s warning? - Qatar has rebuffed Israeli criticism of its mediation, including by Prime Minister Benjamin Netanyahu, for months.
In a statement on Tuesday, the Qatari embassy in Washington also hit out at Congressman Steny Hoyer, the top-ranking @ @ @ @ @ @ @ @ @ @ urged the administration to rethink its relationship with Qatar, and called on Doha to pressure Hamas for a hostage release.
Without naming any individuals, Sheikh Mohammed said Qatar had been the victim of"points-scoring"by"politicians who are trying to conduct election campaigns by slighting the State of Qatar".
Middle East expert James Dorsey said the prime minister ' s statements indicated"Qatar punching back rather than seriously considering giving up mediation", which he called"a key pillar of the country ' s soft power".
He explained while Qatar was primarily targeting"Israel, Benjamin Netanyahu and his backers in the US Congress", there was also"an attempt to pressure the Biden administration to stand up for Qatar".
Gulf expert Andreas Krieg said Qatar had played an"instrumental role"in securing November ' s hostage exchange and the emirate had become"very disgruntled about the way that this isn't recognised by everyone, particularly not in Israel".
But he also said it @ @ @ @ @ @ @ @ @ @ to withdraw from this mediation effort"after it had"monopolised this relationship in a way that no one else can do what Qatar can do". - Why have negotiations stalled? - Israel and Hamas have traded blame over the failure to make headway on a framework put forward by Qatar, the United States and Egypt, which would halt fighting in Gaza for six weeks and see the exchange of about 40 hostages for hundreds of Palestinian prisoners.
Sticking points have included Hamas ' s long-standing demands for a permanent ceasefire and the full withdrawal of Israeli forces from the Gaza Strip, which Israeli officials have repeatedly opposed.
Last week, Israel ' s Mossad spy agency said Hamas ' s rejection of the latest proposal showed the militant group"does not want a humanitarian deal and the return of the hostages". - Would talks end without Qatar? - Hamas ' s political office has been based in Qatar since 2012, at the request of the United States.
Krieg said Qatar had @ @ @ @ @ @ @ @ @ @ had been involved, the"Egyptians don't have the kind of ways in... that Qatar has"."I think the mediation can continue without the Egyptians, but it couldn't continue without the Qataris,"the King ' s College London academic said.
Dorsey said that"if Qatar were to diminish or end its role in the negotiations, then the pressure for it to expel Hamas will increase.""If suddenly it ' s Iran, who do the Americans and the Israelis talk to, to get to Hamas?"Dorsey said.
Sheikh Mohammed announced Qatar ' s reassessment alongside Turkish Foreign Minister Hakan Fidan, who spoke for three hours in Doha on Wednesday with Hamas chief Ismail Haniyeh and other Hamas leaders.
Haniyeh is due to visit Turkey this weekend as a guest of President Recep Tayyip Erdogan."If, for whatever reason, Qatar were forced to bow out or diminish its role in negotiations, Turkey is a candidate,"Dorsey said. @ @ @ @ @ @ @ @ @ @ provided to this website from an external third party provider.
We are not responsible for, and do not control, such external websites, entities, applications or media publishers.
The body of the text is provided on an"as is"and"as available"basis and has not been edited in any way.
Neither we nor our affiliates guarantee the accuracy of or endorse the views or opinions expressed in this article.
Read our full disclaimer policy here.

% matched lemmas: conduit, week-long
\end{textsample}
